\documentclass[12pt]{article}
\usepackage{amsmath,amssymb,amsfonts}
\usepackage[margin=1in]{geometry}

\begin{document}

\section*{Chapter 3, Problem 23}

\textbf{Problem.} Show that if $A^{-1}$ and $B^{-1}$ exist, then $(AB)^{-1} = B^{-1}A^{-1}$.

If $A$ and $B$ are similar, prove that $A^n$ and $B^n$ are similar.

\bigskip
\textbf{Solution.}

\textbf{Part 1:} $(AB)^{-1} = B^{-1}A^{-1}$.

We verify that $B^{-1}A^{-1}$ is the inverse of $AB$ by checking both products:
\[
(AB)(B^{-1}A^{-1}) = A(BB^{-1})A^{-1} = AIA^{-1} = AA^{-1} = I.
\]
\[
(B^{-1}A^{-1})(AB) = B^{-1}(A^{-1}A)B = B^{-1}IB = B^{-1}B = I.
\]
Since both products yield $I$, we have $(AB)^{-1} = B^{-1}A^{-1}$. \hfill $\checkmark$

\bigskip
\textbf{Part 2:} If $A$ and $B$ are similar, then $A^n$ and $B^n$ are similar.

Since $A$ and $B$ are similar, there exists an invertible matrix $P$ such that $B = P^{-1}AP$. Then:
\[
B^n = (P^{-1}AP)^n = (P^{-1}AP)(P^{-1}AP)\cdots(P^{-1}AP).
\]
The intermediate factors telescope: $P(P^{-1}) = I$, so
\[
B^n = P^{-1}A^n P.
\]
This shows $A^n$ and $B^n$ are related by the same similarity transformation $P$, hence they are similar. \hfill $\blacksquare$

\end{document}
